\cotefile{157-2}
\scanpage{67}
\mspage{11}

\uline{Dém} Je dis qu'il suffit de prouver, pour les trois foncteur\unsure{s} $p,q,r$, et pour les foncteurs en sens inverse
\diagram{\begin{tikzcd}[ampersand replacement=\&,column sep=large]
\underline{\mathrm{Fib}\,\mathrm{sc}}_{\unsure{0|o}}\,S
 \& \underline{\mathrm{Parf}}\,S
 \& \underline{\mathrm{Cofib}}\,\underline{\mathrm{sc}}_0\,S \\
 \& \mathrm{Cat}_{/S}\arrow[ul,"\underline{\Phi}_0"']\arrow[u,"C"]\arrow[ur,"\Psi_0"'] \&
\end{tikzcd}}
, que les six composés
\[
p\underline{\Phi}_0,\underline{\Phi}_0p,\quad q\Psi_0,\Psi_0q,\quad rC,Cr
\]
sont reliés aux foncteurs identiques (dans la catégorie pertinente) par une flèche qui est dans le localisateur \deleted{\unsure{envisagé}} correspondant. Pour \unsure{ch. un|chacun} des trois couples de foncteurs \inserted{$(p,\underline{\Phi}_0),(q,\Psi_0),(r,C)$}, cela prouvera déjà qu'en passant aux localisés, \deleted{\unsure{que}} ils induisent des équivalences quasi-inverses l'une de l'autre. Mais on a mieux, à l'aide du lemme formel suivant

\uline{Lemme} Soient $(C,W),(C',W')$ \deleted{$\unsure{(C'',W'')|(C,W)}$} \unsure{deux} catégories avec localisateurs $W\subset\operatorname{Fl}C,W'\subset\operatorname{Fl}C'$, \deleted{$\unsure{W''\subset\operatorname{Fl}C''}$}. Soit
\[ p:C\unsure{\longrightarrow|\text{---}}C' \]
un foncteur compatible avec les localisateurs, et supposons qu'il existe un foncteur $\unsure{\varphi|q}:C'\unsure{\text{---}|\longrightarrow}C$ également compatible, et tel que \deleted{\unsure{l'on}} l'on ait des flèches (à direction non précisée)
\[ \varphi p\overset{u}{\text{---}}\mathrm{id}_C,\qquad p\varphi\overset{v}{\text{---}}\mathrm{id}_{C'} \]
qui soient dans $W$ resp. $W'$ ``\unsure{argument} par argument''.

\deleted{\unsure{Alors, pour tout couple} $(A,\Sigma),\ \Sigma\subset\operatorname{Fl}(A)$, \unsure{un} \unsure{troisième} et \unsure{tout} foncteur $f:A\longrightarrow C'$}
\deleted{\unsure{Soit d'autre part}}

% archive sheet number: 66
\scanpage{68}
\mspage{12}

\deleted{\deleted{\unsure{tel qu'il y|tel qu'il est}} compatible avec les localisateurs, et tel qu'il existe un foncteur $g:C\unsure{\text{---}|\longrightarrow}A$ compatible avec les localisateurs \deleted{\unsure{et tel que}} avec $fg=p$, $f$ induit une équivalence \deleted{\unsure{sur}} \unsure{entre} les catégories localisées}
\[ \deleted{A\Sigma^{-1}\xrightarrow{\approx} C'{W'}^{-1}.} \]
\diagram{\deleted{\begin{tikzcd}[ampersand replacement=\&,column sep=small]
 C\arrow[rr,"g"]\arrow[dr,"p"'] \& \& A\arrow[dl,"f"]\\
 \& C'\arrow[ul,dashed,bend left=42,"q"']\arrow[ur,dashed,bend right=42,"f'"'] \&
\end{tikzcd}}}

\deleted{\uline{Dém} On va définir un foncteur en sens inverse de $f$}
\[ \deleted{f':C'\longrightarrow A,} \]
\deleted{savoir $f'=gq$. Je dis que le couple $(f,f')$ satisfait aux \unsure{m\^{e}mes} hypothèses que le couple $(p,q)$. En effet}
\[ \deleted{f'f=gqf} \]

Alors $p$ et $\unsure{\varphi|q}$ induisent des \deleted{\unsure{quasi}} équivalences quasi-inverses l'une de l'autre en passant aux catégories localisées (elles n'ont pas de \unsure{mérite}).

De plus, pour \unsure{tte|tt} sous-catégorie pleine $C'_0$ de $C'$, telle que \deleted{$\unsure{C}$} $p$ se factorise \inserted{(en $p_0:C\unsure{\text{---}|\to}C'_0$)} par $C'_0$, si on désigne par $W'_0$ le localisateur induit \deleted{\unsure{dans}} $C'_0$ par $C'$, \deleted{\unsure{les de}} alors $p_0:C\to C'_0$ et $\varphi_0=\varphi\vert C'_0:C'\unsure{\text{---}|\to}C$ induisent également des équivalences qu. inv. l'une de l'autre sur les localisés.

En effet, $\varphi_0p_0=(\varphi i)p_0=\varphi(ip_0)=\varphi p\overset{u}{\text{---}}\mathrm{id}_C$, (ici $i:C'_0\hookrightarrow C'$ désigne l'inclusion), et $\deleted{\unsure{g}}F\overset{\mathrm{d\acute ef}}{=}p_0\varphi_0=p_0\varphi i$ satisfait à
\[ iF=\underbrace{ip_0}_{p}\varphi i=(p\varphi)i. \]

Pour définir \deleted{$\unsure{i}$} $F\overset{v_0}{\text{---}}\mathrm{id}_{C'_0}$, qui est dans $W'_0$ arg. par argument

% archive sheet number: 67
\scanpage{69}
\mspage{13}

il suffit de définir
\[ \underbrace{iF}_{\substack{\Vert\\(p\unsure{\varphi|q})i}}\overset{\unsure{\cdot}\,\overline{v_0}}{\text{---}}i\deleted{\unsure{\mathrm{id}}} \]
qui soit dans $W'$ argument par argument, puisque $W'_{\unsure{0|o}}$ induit par $W'$ et que $F\mapsto\deleted{\unsure{F}}iF$
\[ \underline{\mathrm{Hom}}(C'_0,C'_0)\unsure{\longrightarrow|\text{---}}\underline{\mathrm{Hom}}(C'_0,C') \]
est \unsure{pl.} fid. ($i$ l'étant). Or il suffit de prendre $\overline{v_0}=v*i:(p\varphi)\circ i\ \text{---}\ \mathrm{id}_{C'}\circ i$.

Pour résumer, il faut montrer que le coup des couples de foncteurs quasi-inverses marche dans les cinq cas suivants \unsure{:|6}
\diagram{\begin{tikzcd}[ampersand replacement=\&,column sep=small,row sep=huge]
\underline{\mathrm{Fib}\,\mathrm{sc}}\arrow[drr,shift left]
 \&\underline{\mathrm{Fib}\,\mathrm{sc}}_0\arrow[l,hook']\arrow[dr,shift left]
 \&\underline{\mathrm{Parf}}\arrow[d,shift left]
 \&\underline{\mathrm{Cofib}\,\mathrm{sc}}_0\arrow[dl,shift right]
 \&\underline{\mathrm{Cofib}}\arrow[dll,shift right]\\
 \& \&\mathrm{Cat}_{/S}\arrow[ull,shift left,"\underline{\Phi}"']\arrow[ul,shift left,"\underline{\Phi}_0"']\arrow[u,shift left,"C"]\arrow[ur,shift right,"\Psi_0"']\arrow[urr,shift right,"\Psi"']\&\&
\end{tikzcd}}

D'ailleurs par dualité, on est ramené aux \inserted{trois} cas de $\underline{\mathrm{Fib}\,\mathrm{sc}}$, $\underline{\mathrm{Fib}\,\mathrm{sc}}_0$, $\underline{\mathrm{Parf}}$. (Peut être le cas $\underline{\mathrm{Fib}\,\mathrm{Sc}}_0$, avec $\underline{\Phi}_0$ défini en termes de $\underline{\Phi}$ et de $C$, se ramènerait-il aux deux cas qui l'entourent, de $\underline{\mathrm{Fib}\,\mathrm{sc}}$ et $\underline{\mathrm{Parf}}$.)

\centerline{$\unsure{\backslash}$}

Je vois que je dois rectifier le tir, tout d'abord parce que j'ai oublié que les foncteurs

% archive sheet number: 68
\scanpage{70}
\mspage{14}

$\underline{\Phi},\Psi$ que j'avais en vue ne \deleted{\unsure{transforment}} \inserted{sont} \uline{pas} \deleted{\unsure{équivalences} des \unsure{$S$-morphismes}} compatibles avec les localiseurs --- une flèche \inserted{$f:X\unsure{\text{---}|\to}Y$} de $\mathrm{Cat}_{/S}$ \deleted{\unsure{qui est}} telle que la flèche correspondante de $\mathrm{Cat}$ soit dans $W$, ne va pas induire une flèche dans $\underline{\mathrm{Fib}\,\mathrm{sc}}\ S$ ou $\underline{\mathrm{Cofib}\,\mathrm{sc}.}\ S$ qui soit dans $W$ sur chaque fibre.

Ce n'est pas grave, car on peut \deleted{\unsure{les}} remplacer ces foncteurs par $\underline{\Phi}'=\alpha'\underline{\Phi}_{\unsure{0|o}}$, $\Psi'=\beta'\Psi_0$, qui sauf erreur doivent marcher. Mais en vue du lemme précédent, il y a une perplexité, car j'ai l'air d'avoir prouvé que si $\unsure{\mathcal F|\mathscr F}$ est une sous-catégorie pleine de $\mathrm{Cat}_{/S}$ \unsure{\{} dont je suppose seulement qu'elle contient \deleted{\unsure{soit}} l'image essentielle d'une des trois catégories $\underline{\mathrm{Fib}}_0\,S$, $\underline{\mathrm{Cofib}}_0\,\{S\}$, ou $\underline{\mathrm{Parf}}/S$, alors l'inclusion dans $\mathrm{Cat}_{/S}$ induit une équivalence de catégories localisées, en prenant dans $\mathcal F$ comme localisateur celui induit par $\mathrm{Cat}_{/S}$, \unsure{i.e.}
\diagram{\begin{tikzcd}[ampersand replacement=\&,column sep=tiny,row sep=tiny]
 f\colon X\arrow[rr]\arrow[dr,no head]\&\&Y\arrow[dl,no head]\\
 \& S\&
\end{tikzcd}}
est dans $W_{\mathcal F}$ ssi en tant que flèche de $\mathrm{Cat}$, $f$ est dans $W$. Or ceci coïncide bien \deleted{\unsure{au}} \unsure{(via} un théorème non trivial\unsure{,|?} sur lequel je \unsure{devrai} \unsure{revoir|revenir} dans le cadre d'un localisateur fondamental général $W$) avec \unsure{les ensembles|l'ensemble} de flèches de localisation ``fibre par fibre'', dans les cas de $\mathcal F=\underline{\mathrm{Parf}}\ S,\ \underline{\mathrm{Fib}}\,\underline{\mathrm{hot}}\ S,\ \underline{\mathrm{Fib}}_0\,S,\ \underline{\mathrm{Cofib}}_0\,S$.

% archive sheet number: 69
\scanpage{71}
\mspage{15}

Par contre, ce n'est nullement le cas en prenant $\unsure{\mathcal F|\mathscr F}=\underline{\mathrm{Fib}}\,S,\underline{\mathrm{Liss}}\,S$ ou $\underline{\mathrm{Cofib}}\,S,\underline{\mathrm{Prop}}\,S$. Ayant choisi dans ces catégories comme localisateurs \unsure{la notion|les notions} de $W$-équivalences ``fibre par fibre'', j'ai donc l'air de dire que je trouve pour ces quatre catégories les \uline{mêmes} \unsure{localisées|localisés} (à équivalence près) que si je prenais le localisateur beaucoup plus grossier hérité par $\mathrm{Cat}_{/S}$. Cela semble difficile à croire, et il faudra que je fasse très attention en démontrant les énoncés-clef.

Pour voir que les flèches en traits pleins du grand diagramme p.~6 sont toutes compatibles avec les localis(at)eurs, donc aussi les flèches composées au moyen de celles-ci, j'ai à utiliser déjà ceci (pour \inserted{le cas} \deleted{\unsure{de}}des flèches $\underline{\mathrm{Liss}}\,S\unsure{\text{---}|\to}\mathrm{Cat}_{/S},\ \underline{\mathrm{Prop}}_{/S}\to\mathrm{Cat}_{/S}$)\,:

\textcircled{1} Si $X,Y$ lisses \inserted{ou propres} sur $S$, et $f:X\unsure{\text{---}|\to}Y$ un $S$-\deleted{\unsure{fonc}}\unsure{foncteur}, si $f_s:X_s\to Y_s$ est dans $W$ pour \unsure{tout} $s\in S$, alors $f\in W$.
% slanted
\margin{OK si on suppose L 3 gauche et L 4}

Cela renforce considérablement la condition L 5 sur $W$, où on supposait \deleted{\unsure{que}} que $f$ est un foncteur \uline{cartésien} entre catégories \uline{fibrées} sur $S$.

(À \unsure{dégager|déduire}, pour les différentes propriétés dont on a besoin\unsure{$\ne$|,} pour $W$, à quelles propriétés \unsure{standart|standard} on peut les ramener. Tout sera OK si $W=W_\infty$.)

% archive sheet number: 70
\scanpage{72}
\mspage{16}

I) \uline{Définition de} $\underline{\Phi},\Psi$.

On va calquer la construction \inserted{canonique} de $\underline{\mathrm{Ch}}(X)\to X\times X$, \inserted{(ou $\underline{\mathrm{Ch}}_{\unsure{\infty|w}}(X)\unsure{\text{---}|\to}X\times X$\unsure{,})} factorisant $X\xrightarrow{\mathrm{diag}}X\times X$, et la \deleted{\unsure{constr.}} factorisation canonique correspondante d'un $f:X'\unsure{\text{---}|\to}X$ en $X'\xrightarrow{\deleted{\unsure{\sim}}\,i}\overline{X}'\xrightarrow{\bar f}Y$. On remplacera simplement $\underline{\mathrm{Ch}}(X)$ par la sous-catégorie $\underline{\mathrm{Ch}}_{\deleted{\unsure{\Delta}}+}(X)$, laquelle n'est autre que la catégorie des flèches de $X$
\[ \underline{\mathrm{Fl}}(X)=\deleted{\unsure{\mathrm{Hom}}}\,\underline{\mathrm{Hom}}(\Delta_1,X), \]
muni des foncteurs
\[ X\xrightarrow{i_X}\underline{\mathrm{Fl}}(X)\xrightarrow{\inserted{p_X=(\sigma_X,\beta_X)}} X\times X, \]
où $i_X$ est la flèche $x\longmapsto\mathrm{id}_x$, et \deleted{$\unsure{i_X(x)=}$} $\sigma_X,\beta_X$ sont les foncteurs source et but. La fibre de $p_X$ \inserted{en $(a,b)\in\operatorname{Ob}(X\times X)$} est la catégorie discrète définie par l'ens $\mathrm{Hom}(a,b)$, \unsure{ça} dépend \unsure{grandement} de $(a,b)$ et pas seulement de sa composante connexe dans $X\times X$ --- donc $p_X$ n'a rien d'une Hot-fibration. Il n'est en général ni propre ni lisse\unsure{,} ni rien, par contre les deux flèches composées avec $\mathrm{pr}_1$ et $\mathrm{pr}_2$
\diagram{\begin{tikzcd}[ampersand replacement=\&,column sep=small]
\&\underline{\mathrm{Fl}}(X)\arrow[dl,"\sigma_X"']\arrow[dr,"\beta_X"]\&\\
X\&\& X
\end{tikzcd}}
sont sympathiques. on a en effet

% archive sheet number: 71
\scanpage{73}
\mspage{17}

\uline{Lemme} \inserted{On a} $\sigma_X^{-1}(\{x\})\simeq{}_{x\backslash}X,\quad \beta_X^{-1}(\{y\})\simeq X_{/x}$. \deleted{\unsure{Le}}

Le foncteur $\sigma_X$ est \inserted{Cat-}fibrant à fibres \deleted{\unsure{asphé}} contractibles (donc asphériques, donc $W$-asphériques), le foncteur $\beta_X$ cofibrant à fibres contractibles (donc asphériques donc $W$-asphériques). Ce sont donc \inserted{aussi} des Hot-équivalences a fortiori des $W$-équivalences.
% slanted
\margin{En fait, ce sont même des \deleted{\unsure{équiv. d'}} \unsure{homotopismes}, cf. \unsure{ci-après|ci-dessous} (un argument plus intrinsèque \unsure{s'invite|s'insère} \unsure{naturellement dans|en particularisant} le cas de $X'=X$)}

\uline{Corollaire 1} Le foncteur $i_X:X\unsure{\text{---}|\to}\underline{\mathrm{Fl}}(X)$ est une Hot-équivalence (puisque \unsure{sa composée|son composé} avec la Hot-équivalence \inserted{$\sigma_X$} est l'identité).

Voici la construction générale de factorisation de $f:X'\longrightarrow X$, correspondant au choix précédent d'une ``catégorie des chemins'' $\unsure{\mathcal C|C}(X)=\underline{\mathrm{Fl}}(X)$
\diagram{\begin{tikzcd}[ampersand replacement=\&,row sep=small,column sep=normal]
\overline{X}'\arrow[dr]\arrow[drr,"\bar f"]\arrow[dd,"\pi_f" description,double] \&\&\\
\&\underline{\mathrm{Fl}}(X)\arrow[r,"\beta_X"']\arrow[d,"\sigma_X"]\&X\\
X'\arrow[uu,dashed,bend left=27,"i_f"]\arrow[r,"f"']\&\deleted{\unsure{X}}\unsure{X}\&
\arrow[from=1-1,to=3-2,phantom,"\mathrm{cart}" description]
\end{tikzcd}}
% slanted
\margin{\uline{NB} $\pi_f\cdot i_f=\mathrm{id}$, et $i_f\pi_f\longrightarrow\mathrm{id}$, donc $i_f,\pi_f$ sont des homotopismes, \deleted{\unsure{donc}} inverses l'un de l'autre dans $\mathrm{Cat}\,\mathrm{hot}$. Cela précise cor.~3 page 18.}
où $i_f$ est déduit de la section $i$ de $\underline{\mathrm{Fl}}(X)$ sur $X$. Ainsi, $\overline{X}'$ est la catégorie des couples
\[ (x',u)\mid x'\in\operatorname{Ob}X',\quad \underset{\substack{\text{\tikz[baseline=-.4ex]\node[rotate=-90,inner sep=0pt] {$\scriptstyle\in$};}\\\mathrm{Fl}(X)}}{u}:f(x')\longrightarrow y \]
le foncteur $\bar f$ est défini par
\[ \bar f(x',u)=\mathrm{but}(u)=y \]
et $i_f$ par
\[ i_f(x')=(x',\mathrm{id}_{f(x')}:f(x')\overset{\mathrm{id}}{\text{---}} f(x')) \]
donc on a bien $f(x')=\bar f\deleted{\unsure{\varphi}}\,i_f(x')$.

% archive sheet number: 72
\scanpage{74}
\mspage{18}

\uline{Cor. 2} Le foncteur
\[ \bar f:\overline{X}'\unsure{\longrightarrow|\text{---}}X \]
a comme fibres les catégories $X'_{/x}$, il est cofibrant, les \deleted{\unsure{applications}} \inserted{foncteurs} de changement de base sont les foncteurs $X'_{/x}\unsure{\text{---}|\to}X'_{/y}$ associés aux \deleted{\unsure{foncteurs}} \inserted{flèches} $x\unsure{\text{---}|\to}y$ de $X$.

En fait, on voit que cette catégorie est non seulement cofibrée, mais canoniquement \uline{scindée}. \deleted{\unsure{On a donc}}
% slanted
\margin{\uline{Cor. 3}\,: La flèche $i_f:X'\unsure{\text{---}|\to}\overline{X}'$ est un \deleted{$\unsure{+|x}$} $X$-morphisme (\unsure{ici} $X'$ sur $X$ par $f$, $\overline{X}'$ par $\bar f$)\,: \unsure{i.e.} $f=\bar f\circ i_f$, et c'est une Hot-équivalence (Car $\pi_f:\overline{X}'\to X'$, \unsure{est} Cat-fibrant à fibres contractibles, donc une Hot-équivalence)}

Appliquant ceci à des $f:X\unsure{\text{---}|\to}S$, on trouve une construction
\[ \mathrm{Cat}_{/S}\xrightarrow{\deleted{\unsure{\underline{\Phi}}}\Psi_{\deleted{\unsure{\mathrm{sc}}}}}\underline{\mathrm{Cofib}\,\mathrm{sc}}\ S \]
c'est bel et bien fonctoriel en $X$ variant dans $\mathrm{Cat}_{/S}$. Si \unsure{on a}

\deleted{On définit $\underline{\Phi}$ par composition}
\diagram{\deleted{\begin{tikzcd}[ampersand replacement=\&,column sep=large]
\mathrm{Cat}_{/S}\arrow[r,"\underline{\Phi}_{\mathrm{sc}}"]\arrow[dr,"\underline{\Phi}"']\&\underline{\mathrm{Cofib}\,\mathrm{sc}}\ S\arrow[d]\\
\&\underline{\mathrm{Cofib}}\ S
\end{tikzcd}}}
\[ \unsure{\varphi|\psi}:X\xrightarrow{\deleted{\unsure{f}}}Y \]
dans $\mathrm{Cat}_{/S}$, alors $\varphi$ induit des
\[ \varphi_{/s}:X_{/s}\longrightarrow Y_{/s} \]
pour $s\in S$, et ceci de façon compatible avec les

% archive sheet number: 73
\scanpage{75}
\mspage{19}

foncteurs ``\unsure{cochgt.|cochgt} de base'' associés aux $u:s\unsure{\text{---}|\to}t$ dans $S$
\diagram{\begin{tikzcd}[ampersand replacement=\&]
 X_{/s}\arrow[r,"u_*^X"]\arrow[d,"\unsure{\varphi|\psi}_{/s}"'] \& X_{/t}\arrow[d,"\varphi_{/t}"]\\
 Y_{/s}\arrow[r,"u_{\unsure{*|x}}^Y"'] \& Y_{/t}
\end{tikzcd}}

Cela définit donc un morphisme de catégories scindées sur $S$
\[ \deleted{\underline{\Phi}}\Psi(X)\longrightarrow\deleted{\underline{\Phi}}\Psi(Y), \]
c'est lui !

On a \deleted{des} \inserted{un} morphisme\deleted{\unsure{s}}, fonctoriel\unsure{,} en $X$
\diagram{\begin{tikzcd}[ampersand replacement=\&]
 X\arrow[r,shift left,"i_X"] \& \deleted{\underline{\Phi}}\Psi(X)\arrow[l,shift left,"\deleted{\pi_X}",description,draw=none,"\deleted{\longleftarrow}"']
\end{tikzcd}}
\inserted{dans $\mathrm{Cat}_{/X}$} \deleted{$\unsure{\pi_Xi_X=\mathrm{id}_X}$} et c'est une Hot-équivalence, donc $\in W$.

\deleted{fonctoriel en $X$ --- qui est dans $W_\infty$ avec $\pi_X$ Cat-fibrant à fibres contractibles, donc une Hot-équivalence, donc une $W$-équivalence. (NB dans le diagramme p.~17, $\pi_f$ est déduit de $\sigma_X$, qui est fibrant à fibres contractibles, par un chgt. de base $f$. Mais Attention, le $X$ \unsure{de page|des pages} 17 est devenu $S$, le $X'$ est devenu $X$.) Donc $i_X$ est également une Hot-équivalence.}

Sur les fibres, les morphisme\unsure{s} $i_X$, \deleted{$\unsure{\pi_X}$} induis\unsure{ent|ent-ils} \deleted{\unsure{des}}
\diagram{\begin{tikzcd}[ampersand replacement=\&]
X_s\arrow[r,shift left,"i_{X,s}"]\& X_{/s}\arrow[l,shift left,draw=none,"\deleted{\pi_{X,s}}"',"\deleted{\longleftarrow}"]
\end{tikzcd}}
, \deleted{le foncteur $i_{X,s}$ est} l'inclusion canonique, \deleted{le foncteur}

% archive sheet number: 74
\scanpage{76}
\mspage{20}

\uline{Corollaire 1} Si $X\xrightarrow{f}Y$ de $\mathrm{Cat}_{/S}$ est dans $W$, alors $\deleted{\unsure{\Phi}}\Psi(X)\xrightarrow{\Psi(f)}\Psi(Y)$ aussi. Mais on se gardera de croire que c'est dans $W_S$ \unsure{du|au} sens de $\underline{\mathrm{Cofib}}\,S$ ou $\underline{\mathrm{Cofib}\,\mathrm{sc}}\,S$.

Car sur les fibres \inserted{de $\Psi(X),\Psi(Y)$} \deleted{c'est} on trouve les flèches
\[ f_{/s}:X_{/s}\unsure{\longrightarrow|\text{---}}Y_{/s} \]
induites par $f$, qui n'ont aucune raison d'être dans $W$. Qu'elles le soient \unsure{signifie|signifierait} que $f:X\unsure{\text{---}|\to}Y$ est \unsure{dans} $W$ ``localement au dessus de $S$'', ce qui \inserted{(pour les $W$ qui nous intéressent)} est beaucoup plus fort que d'être dans $W$. (Si $Y=S$, cela signifie en fait que $p:X\to S$ est $W$-asphérique, beaucoup plus fort que \deleted{\unsure{le fait}} $p\in W$.)
% slanted
\margin{cf \textcircled{1} \unsure{$\blacksquare$} page 15.}

Donc \deleted{$\unsure{\Psi}$} pour prouver que
\[ \unsure{\text{``}}\,\underline{\mathrm{Cofib}}\,S\to\mathrm{Cat}_{/S}\quad\text{ou}\quad \underline{\mathrm{Cofib}\,\mathrm{sc}}\,S\unsure{\text{---}|\to}\mathrm{Cat}_{/S} \]
induisent des équivalences pour les catégories localisées, le foncteur $\Psi$ directement ne \deleted{\unsure{nous}} sera pas utilisable. Cependant\,:

\uline{\unsure{Corollaire} 2} Supposons $X,Y$ dans $\underline{\mathrm{Prop}}\,S$ \inserted{ou dans $\mathrm{Fib\,hot}\,S$}, et considérons $\Psi(f):\Psi(X)\to\Psi(Y)$, pour que ce soit une $W$-équivalence dans $\underline{\mathrm{Cofib}\,\mathrm{sc}}\,S$ \deleted{(ou dans $\underline{\mathrm{Cofib}}\,S$, ce qui revient au même),}

% archive sheet number: 75
\scanpage{77}
\mspage{21}

i.e. une $W$-équivalence fibre par fibre, il faut \unsure{s.|?.} que $f$ soit une $W$ équivalence fibre par fibre (et cela doit impliquer \inserted{cf \textcircled{1} p.~15} que $f$ est une $W$-équivalence)

En effet, dans le diagramme commutatif
\diagram{\begin{tikzcd}[ampersand replacement=\&]
X_s\arrow[r,"f_s"]\arrow[d]\&Y_s\arrow[d]\\
X_{/s}\arrow[r,"f_{/s}"']\&Y_{/s}
\end{tikzcd}}
les flèches verticales sont dans $W$, car \deleted{\unsure{f est|sont}} \inserted{$X,Y$} propres sur $S$ \inserted{(ou $X,Y$ dans $\underline{\mathrm{Fib}\,\mathrm{hot}}\,S$)}, donc $f_s$ est dans $W$ ssi $f_{/s}$ l'est.
% margin note written sideways
\margin{\uline{Cor. 3} Si $X\in\operatorname{Ob}\underline{\mathrm{Fib}\,\mathrm{hot}}\,S$, \unsure{$\blacksquare$} alors $\Psi(X)$ est dans $\underline{\mathrm{Fib}\,\mathrm{sc}}_{\unsure{0|o}}\,S$. \deleted{\unsure{On regarde}} Cela signifie en effet que les $X_{/s}\unsure{\text{---}|\to}X_{/t}$ sont dans $W$, or l'hyp. sur $X$ \unsure{\{} \unsure{indique|implique} que c'est des $\{W_\infty\}$\unsure{,} OK.}
% margin note written sideways
\margin{\diagram{\deleted{\begin{tikzcd}[ampersand replacement=\&,column sep=small,row sep=small]
 X_s\arrow[r]\arrow[d]\&X_t\arrow[d]\\
 X_{/s}\arrow[r]\& X_{/t}
\end{tikzcd}}}}

$\underline{\Phi}$ se définit de façon duale à $\Psi$, ce qui revient à échanger les rôles de $\sigma_X,\beta_X$ dans le diagramme page 17.

II) Définition de \deleted{\unsure{Parf}} $C:\mathrm{Cat}_{/S}\unsure{\text{---}|\to}\underline{\mathrm{Parf}}\ \deleted{\unsure{X}}S$.

Cette fois, c'est la factorisation de \deleted{\unsure{toute}} flèches $\deleted{\unsure{f}}p:X\longrightarrow S$ dans $\mathrm{Cat}$, en
\[ X\xrightarrow{j_{X/S}}\overline{X}\xrightarrow{\bar p} S, \]
où $i_X$ est une Hot-équivalence (donc dans $W$), et $\bar p$ parfait. Cela se construit formellement comme tantôt, en utilisant $\underline{\mathrm{Ch}}_\infty(\deleted{\unsure{X}}S)$ au lieu de $\underline{\mathrm{Fl}}(\deleted{\unsure{X}}S)$, et la factorisation
\[ \deleted{\unsure{X}}S\xrightarrow{\deleted{\unsure{i}}j_S}\underline{\mathrm{Ch}}_\infty(\deleted{\unsure{X}}S)\xrightarrow{p_S=(\sigma_S,\beta_S)}\deleted{\unsure{X}}S\times\deleted{\unsure{X}}S \]
de $\mathrm{diag}_S:S\longrightarrow S\times S$. La situation est plus jolie, car $p_S$ est déjà ``parfait'', ce qui \unsure{impliquera|implique} formellement que le $\bar p$ qui s'en déduit, pour factoriser $p:X\longrightarrow S$, est lui-même parfait.

% archive sheet number: 76
\scanpage{78}
\mspage{22}

\deleted{\unsure{Comme}} \deleted{Si on a une $W$-équivalence}

Ici encore, la construction est fonctorielle, on trouve le foncteur
\[ C:\mathrm{Cat}_{/S}\longrightarrow\underline{\mathrm{Parf}}\,S \]
qu'il faut comparer au foncteur en sens inverse
\[ r:\mathrm{Cat}_{/S}\hookleftarrow\underline{\mathrm{Parf}}\,S. \]

Tout d'abord, vérifions que $C$, \deleted{\unsure{tout}} comme $r$, est compatible avec les \deleted{flèches} localiseurs.

Soit $f:X\to Y$ dans $\mathrm{Cat}_{/S}$, supposons $f$ est dans $W$, \deleted{alors} \inserted{prouvons $C(\unsure{s|S|f})$ est dans $W_S$.} On a commutativité \unsure{du|des}
\diagram{\begin{tikzcd}[ampersand replacement=\&]
\overset{\deleted{\unsure{C}}}{X}\arrow[r,"f"]\arrow[d,"j_X"']\&Y\arrow[d,"j_Y"]\\
CX\arrow[r,"Cf"']\&CY
\end{tikzcd}}
et $j_X,j_Y$ sont dans $W$, donc $f$ est dans $W$ ssi $Cf$ l'est. \deleted{\unsure{que}} Mais dans le cas $W=W_\infty$, on sait que cela équivaut \unsure{à|:} ce que $Cf$ soit dans $W_\infty$ fibre par fibre. Admettant \textcircled{1} p.~15 pour $W$, il reste à s'assurer \deleted{que} d'une réciproque \unsure{dans|des|sur} $\underline{\mathrm{Parf}}$, i.e.

\textcircled{2} Soit
\diagram{\begin{tikzcd}[ampersand replacement=\&,column sep=tiny,row sep=tiny]
X\arrow[rr,"f"]\arrow[dr,no head]\&\&Y\arrow[dl,no head]\\
\&S\&
\end{tikzcd}}
dans $\underline{\mathrm{Parf}}_{/S}$ (ou seulement dans $\underline{\mathrm{Fib}\,\mathrm{hot}}\,S$). Si $f$ est dans $W$, alors les $f_s:X_s\unsure{\text{---}|\to}Y_s$ aussi.

% archive sheet number: 77
\scanpage{79}
\mspage{23}

Essayons de le prouver, en calquant la dém. dans le cas $C=C_\infty$ (\uline{XI}, th.~2). On peut supposer $S$ \uline{connexe}, en \deleted{supposant} utilisant $\mathrm{L}\,\unsure{\mathrm{a}|0}$ (i.e. $W\subset W_{\unsure{0|o}}$) et en admettant le
% slanted
\margin{ne sert à rien !}

$\blacklozenge$ \uline{Lemme}. Soit $f:X\to Y$ dans $\mathrm{Cat}$. Alors $f\in W$ ssi $\pi_0(f)$ est bijectif, et pour \unsure{tte} composante \unsure{conn.|connexe} $X_i$ de $X$, \unsure{\{|$f$} si $Y_i$ est la comp. conn. corr. de $Y$, $f_i:X_i\to Y_i$ induit par $f$ est dans $W$.

Utilisons alors l'existence de
\[ S'\longrightarrow S \]
Hot-fibration avec $S'$ Hot-asphérique. \unsure{OPS} = $S'\unsure{\text{---}|\to}S$ est parfait. Supposons qu'on sache qu'un chgt. de base par de telles fibrations parfaites transforme \deleted{chgt. de} $W$-équivalence en itou. (Ça doit être vrai pour les $W_i$, mais on a vu que c'était faux pour l'équivalence \unsure{ho-homologique}, cf I p.~87 ff) Alors
% slanted
\margin{\uline{NB} Il suffit d'avoir une Hot-fibration $S'\unsure{\text{---}|\to}S_0\subset S$, où $S_0$ est la comp. conn. de $S$ contenant l'élément $s\in\operatorname{Ob}S$ envisagé.}
$X'\xrightarrow{f'}Y'$ sur $S'$ est une $W$-équivalence. Or pour $s'\in S'$, comme $\{s'\}\longrightarrow S'$ est une Hot-équivalence ($S'$ Hot-asphérique), et $X',Y'$ étant Hot-fibrés sur $S$, il s'ensuit que $X'_{s'}\to X',\ Y'_{s'}\unsure{\text{---}|\to}Y'$ sont des Hot-équivalences, donc aussi $X_{s'}\overset{f'_{s'}}{\text{---}}Y_{s'}$, comme on voit sur le carré comm.
\diagram{\begin{tikzcd}[ampersand replacement=\&]
X'_{s'}\arrow[r,"f'_{s'}"]\arrow[d,"\sim"']\& Y'_{s'}\arrow[d,"\sim"]\\
X'\arrow[r,"f'"']\&Y'
\end{tikzcd}}

% archive sheet number: 78
\scanpage{80}
\mspage{24}

Comme $S'\unsure{\text{---}|\to}S$ est surjectif, cela montre que les $f_s:X_s\unsure{\text{---}|\to}Y_s$ sont dans $W$, qed.

Ainsi, (2) sera conséquence de

\textcircled{$2_0$} Soit $p:\deleted{\unsure{X'}}Y'\to\deleted{\unsure{X}}Y$ une Hot-fibration (ou \inserted{même} un foncteur parfait, si ça peut marcher \ldots). Soit $f:X\to Y$ une $S$-flèche (dans $\mathrm{Cat}_{/S}$), \unsure{déduis|déduisons|obtenu} $f':X'\unsure{\text{---}|\to}Y'$ par chgt. de base. Si $f\in W$, $f'\in W'$\,:
\diagram{\begin{tikzcd}[ampersand replacement=\&,column sep=small,row sep=small]
X'\arrow[r,"q"]\arrow[d,"f'"']\&X\arrow[d,"f"]\\
Y'\arrow[r,"p"']\&Y
\arrow[from=1-1,to=2-2,phantom,"\mathrm{cart}" description]
\end{tikzcd}}
\[ f\in W\ \text{et}\ p\ \text{parfait}\ \deleted{\unsure{\text{sur}}}\ \Longrightarrow f'\in W \]

En tous cas, il nous faut absolument \textcircled{2} pour pouvoir conclure que $C$ est compatible avec les localiseurs.

Considérons $rC:\mathrm{Cat}_{/S}\unsure{\longrightarrow|\text{---}}\mathrm{Cat}_{/S}$ on a une flèche
\[
\left\{\begin{aligned}
&\mathrm{id}_{\mathrm{Cat}_{/S}}\xrightarrow{\deleted{\unsure{i}}j}rC\\
&\deleted{\unsure{i}}j_X:X\longrightarrow CX\qquad \substack{\text{la flèche canonique}\\ i_{X/S}\ (\text{p.~21})},
\end{aligned}\right.
\]
et $j_X$ est tjrs dans $W$.

Considérons $Cr:\underline{\mathrm{Parf}}\,S\unsure{\text{---}|\to}\underline{\mathrm{Parf}}\,S$, i.e appliquons la construction $C$ à un $X$ parfait sur $S$. Alors on a encore
\[
\left\{\begin{aligned}
&\mathrm{id}_{\underline{\mathrm{Parf}}\,S}\xrightarrow{j'}Cr\\
&j'_X=j_{\unsure{X|r}}(X):\overset{\deleted{\unsure{r(X)}}}{X=r(X)}\xrightarrow{j_X}C(X).
\end{aligned}\right.
\]

Ledit morphisme est dans $W$, donc aussi une $W$-équiv fibre par fibre, puisque \unsure{sources|source} et \unsure{buts|but} \deleted{sont} parfaits, donc des Hot-fibrés sur $S$.

% archive sheet number: 79
\scanpage{81}
\mspage{25}

Cela prouve déjà que les foncteurs $r$ et $C$ induisent des équivalences \unsure{quasi-inverses} l'une de l'autre sur les catégories localisées. Par le lemme p.~11, cela montre que pour tte sous-catégorie pleine \inserted{$\unsure{\mathcal F|\mathscr F|F}$} de $\mathrm{Cat}_{/S}$ contenant $\underline{\mathrm{Parf}}\,S$, l'inclusion de $\mathcal F$ dans $\mathrm{Cat}_{/S}$ induit une équivalence sur les localisés, en prenant comme localiseur dans $\mathcal F$ le localiseur ``grossier'' induit par $\mathrm{Cat}_{/S}$. Ceci s'applique notamment aux sous-catégories pleines
\[
\underline{\mathrm{Liss}}\,S,\quad\underline{\mathrm{Fib\,hot}}\,S,\quad\underline{\mathrm{Prop}}\,S,
\]
% vertical stroke beside the following four manuscript lines
\leavevmode\llap{\smash{\rule[-2.3\baselineskip]{.4pt}{3\baselineskip}}\quad}mais en prenant dans $\underline{\mathrm{Liss}}\,S$ et $\underline{\mathrm{Prop}}\,S$ \unsure{n'ont|non} pas le localiseur fin ``fibre par fibre'', comme je l'avais annoncé à tort, mais le localiseur grossier induit par $\mathrm{Cat}_{/S}$. (Pour $\underline{\mathrm{Fib\,hot}}\,S$, il n'y a pas de risque d'ambiguïté). Il faudrait que je me \unsure{convainque} que si je prends là les $W$-équivalences fines (dans le cas $W=W_{\infty}$, disons), je ne trouve \uline{pas} les ``bons'' localisés -- i.e. le foncteur\unsure{,|.} des localisés, vers le localisé de $\mathrm{Cat}_{/S}$, n'est pas une équivalence.

Ceci me fait \unsure{supposer} que dans $\underline{\mathrm{Fib}}\,S$ et $\underline{\mathrm{Cofib}}\,S$ également, c'est \deleted{\unsure{le|la}}l'\deleted{\unsure{notion|relation}}\inserted{équivalence} \uline{grossière}

% archive sheet number: 80
\scanpage{82}
\mspage{26}

qu'il convient de prendre. (Dans $\underline{\mathrm{Fib}}_0\,S$ et $\underline{\mathrm{Cofib}}_{\unsure{0}}\,S$, \inserted{ou $\underline{\mathrm{Fibsc}}_0(S)$ et $\mathrm{Cofibsc}_0(S)$,} les deux équivalences sont identiques, puisque ces \unsure{quatre} catégories sont contenues dans $\underline{\mathrm{Fib\,hot}}\,S$.) La suite nous montrera s'il en est bien ainsi.

Par contre, dans $\underline{\mathrm{Fibsc}}\,S$ et $\underline{\mathrm{Cofibsc}}\,S$, il n'y a pas de doute, suivant les principes généraux de construction d'un dérivateur à partir d'une $(\unsure{\mathcal M|M},W)$, qu'il faut prendre l'équivalence fine, fibre par fibre. À ce sujet, il faut faire attention que les foncteurs verticaux du diagramme p.~8
\[
\underline{\mathrm{Fibsc}}\,S\mathbin{\unsure{-|\longrightarrow}}\underline{\mathrm{Fib}}\,S,\qquad
\underline{\mathrm{Cofibsc}}\,S\mathbin{-}\underline{\mathrm{Cofib}}\,S,
\]
quoique \deleted{\unsure{ont}} ess. surjectifs \inserted{et fidèles}, ne sont pas pl. fidèles. Ils le seraient, \inserted{\unsure{et =|=}\ des équiv. de cat.} si dans $\underline{\mathrm{Fib}}\,S$ \deleted{\unsure{et|et Cofib}} $\underline{\mathrm{Cofib}}\,S$ je me bornais aux $S$-foncteurs \uline{cartésiens} resp. \uline{cocartésiens}.

Donc on peut \unsure{interpréter}, à équivalence près, $\underline{\mathrm{Fib}}\,\underline{\mathrm{sc}}\,S$ et $\underline{\mathrm{Cofib}}\,\underline{\mathrm{sc}}\,S$ comme $\underline{\mathrm{Fib}}\,S$ et $\underline{\mathrm{Cofib}}\,S$, où on aurait \unsure{restreint} les morphismes aux seuls $S$-\uline{f}\unsure{oncteurs} \uline{cartésiens}, resp. cocartésiens. Il semblerait donc qu'en restreignant ainsi les flèches, mais en restreignant de façon plus \unsure{draconienne} encore les localiseurs, on trouve des catégories \inserted{localisées} \unsure{coïncidentes|canoniquement|équivalentes}

% archive sheet number: 81
